\chapter{The spectral reading}

\textbf{Purpose:} Read the dependency matrix with every cell rather than its largest, which finds a law the maximum cannot see and then confirms the null harder than the maximum can. \textbf{Scope:} \texttt{src/\allowbreak{}bench/\allowbreak{}bench\_\allowbreak{}sac.\allowbreak{}cu}.

\section{Why a maximum was the wrong reading}

Every strict avalanche reading in the literature and in this work takes the largest of \(131{,}072\) cells and discards the other \(131{,}071\). That is what makes it peak at \(\sqrt{2 \ln N}\) whatever the function does. A structure spread thinly across many cells is invisible to a maximum and plain to a sum.

SHA-256 is built from rotations, and a rotation makes bit \(j\) depend on bit \(j - r\). If any of that survives into the dependency matrix it appears as a function of \((\text{input bit} - \text{output bit}) \bmod 32\), spread over every word pair rather than concentrated anywhere. Folding the cells onto that residue puts \(4096\) of them behind each number instead of one.

The fold is signed rather than absolute, deliberately. Taking magnitudes first would fold noise into a positive bias and manufacture a signal in every class.

\section{A ridge that does not move}

\begin{center}
\begin{tabular}{rrrrl}
\toprule
rounds & residue & sigmas & walked & verdict \\
\midrule
12 & 26 & 15773.15 & --- & structure \\
13 & 26 & 13724.82 & \(+0\) & structure \\
14 & 26 & 11676.61 & \(+0\) & structure \\
15 & 26 & 9627.78 & \(+0\) & structure \\
16 & 26 & 7580.54 & \(+0\) & structure \\
17 & 26 & 5546.15 & \(+0\) & structure \\
18 & 26 & 3551.43 & \(+0\) & structure \\
19 & 26 & 1916.81 & \(+0\) & structure \\
20 & 26 & 789.13 & \(+0\) & structure \\
21 & 26 & 160.61 & \(+0\) & structure \\
22 & 1 & 13.13 & \(+7\) & collapsed \\
23 & 12 & 2.09 & --- & flat \\
32 & 8 & 2.07 & --- & flat \\
64 & 31 & 2.04 & --- & flat \\
\bottomrule
\end{tabular}
\end{center}

\textbf{The ridge does not walk.} Residue twenty-six dominates for ten consecutive rounds. It is a fixed point of the composition rather than something rotating, which answers the question of whether an inverse rotation could unwind it: there is no moving structure to chase.

Not claimed: which rotation residue twenty-six corresponds to. After twelve rounds the effective residue is a composition of many rotations, and attributing it to \(\Sigma_1\)'s \(\mathrm{ROTR}25\) or \(\Sigma_0\)'s \(\mathrm{ROTR}22\) because they sit nearby would be reading a coincidence.

\section{The decay is quantized}

The first four differences are \(2048.33\), \(2048.21\), \(2048.83\) and \(2047.24\). One sigma of the class mean at that sample is \(1.526 \times 10^{-5}\), so \(2048\) sigmas is exactly \(1/32\).

\textbf{The ridge loses one thirty-second of its amplitude per round} --- one bit position of thirty-two going diffuse --- holding to four decimal places for five rounds, then accelerating into collapse at twenty-two. It is a quantized linear decay and not a smooth exponential, which is what a function made of discrete operations should produce and what a continuous model would not predict.

\subsection{Verified at four times the sample}

The law came out of a sample where \(2048\) sigmas happened to equal one thirty-second, and \(2048\) is \(2^{11}\) while the trials were \(2^{18}\). Every clean power of two in that arithmetic came from the chosen sample size, and a law exact at only one sample size is the shape of the measurement rather than a law.

\begin{center}
\begin{tabular}{lrrr}
\toprule
 & \(2^{18}\) trials & \(2^{20}\) trials & ratio \\
\midrule
round 12 & 15773.15 & 31546.74 & 2.0001 \\
round 13 & 13724.82 & 27450.50 & \\
round 14 & 11676.61 & 23354.13 & \\
\midrule
differences & 2048.33 & 4096.24 & \\
 & 2048.21 & 4096.37 & \\
\bottomrule
\end{tabular}
\end{center}

Sigma doubled for four times the trials, which is \(\sqrt{k}\) exactly, while the absolute deviation held at \(4096 \times 2^{-17} = 1/32\). That is the signature of signal, and the precise opposite of the refuted avalanche dip, which held its sigma and shrank its deviation.

\section{The strongest null in this work}

Past round twenty-three the fold reads \(2.09\), \(2.07\) and \(2.04\) against a null peak of \(2.63\), at three depths, with \(4096\) cells behind each number rather than one. That is roughly sixty-four times the sensitivity the maximum had, finding nothing.

The harmonic sweep agrees. The round-twenty-three profile wanders within two sigmas with no visible pattern, and its discrete transform peaks at \(f_{13}\) with magnitude \(1.80\) against a null peak of \(2.35\): no repeating trend, no harmonics, no side lobes.

\section{What the spectral reading settles}

\textbf{Structure exists and is lawful} at rounds twelve to twenty-one: a fixed residue, a quantized decay, verified at two sample sizes.

\textbf{It is gone by round twenty-three}, at sixty-four times the sensitivity of every earlier reading.

\textbf{The unwinding idea is answered negatively} rather than left untried. Nothing rotates, so there is nothing an inverse rotation could follow.
